\documentclass[../../../include/open-logic-section]{subfiles}

\begin{document}

\olfileid{sth}{choice}{justifications}
\olsection{ملاحظات درونی دربارهٔ انتخاب}

پس پرسش کلی‌تر این است که آیا خوش‌ترتیبی، یا انتخاب، یا در واقع
مقایسه‌پذیری همهٔ مجموعه‌ها از نظر اندازه — فرقی نمی‌کند کدام‌یک —
توجیه‌پذیر است.

اینک کوششی برای توجیهی \emph{درونی}. در
\olref[z][story]{sec} چند اصل دربارهٔ سلسله‌مراتب معرفی کردیم. یکی
از آن‌ها ارزش بازگویی دارد:
\begin{enumerate}
	\item[] \stagesacc. برای هر مرحلهٔ $S$ و هر مجموعه‌ای که
	\emph{پیش از} مرحلهٔ $S$ تشکیل شده است، در مرحلهٔ $S$ مجموعه‌ای
	تشکیل می‌شود که اعضایش دقیقاً همان مجموعه‌ها هستند. هیچ چیز دیگری
	در مرحلهٔ~$S$ تشکیل نمی‌شود.
\end{enumerate}
در واقع، نویسندگان بسیاری پیشنهاد کرده‌اند که اصل موضوع انتخاب را
می‌توان با (چیزی مانند) این اصل توجیه کرد. شرح کوتاهی از این رویکرد
ارائه می‌کنیم.

با نتیجهٔ کوچک و ساده‌ای آغاز می‌کنیم که \emph{یک هم‌ارز دیگر} برای
انتخاب عرضه می‌کند:

\begin{thm}[در $\ZF$]\ollabel{choiceset}
انتخاب با اصل زیر هم‌ارز است. اگر !!{element}های $A$ ناتهی و
دوبه‌دو مجزا باشند، آن‌گاه $C$ای وجود دارد به‌طوری‌که برای هر
$x \in A$، مجموعهٔ $C \cap x$ تک‌عضوی است. (چنین $C$ای را
{مجموعهٔ انتخاب} برای $A$ می‌نامیم.)
\end{thm}

اثبات این نتیجه سرراست است و آن را به‌عنوان تمرین به خواننده
واگذار می‌کنیم.

\begin{prob}
\olref[sth][choice][justifications]{choiceset} را اثبات کنید. اگر
دچار مشکل شدید، می‌توانید اثباتی را در
\cite[pp.~242--3]{Potter2004} بیابید.
\end{prob}

نکتهٔ اساسی این است که مجموعهٔ انتخاب برای $A$ همان برد یک تابع
انتخاب برای $A$ است. بنابراین، برای توجیه انتخاب کافی است بکوشیم
صورت‌بندی هم‌ارز آن را بر حسب وجود مجموعه‌های انتخاب توجیه کنیم.
اکنون دقیقاً همین کار را خواهیم کرد.

فرض کنید !!{element}های $A$ ناتهی و دوبه‌دو مجزا باشند. بنا بر
\stageshier{} (نگاه کنید به \olref[z][story]{sec})، $A$ در مرحله‌ای
مانند~$S$ تشکیل می‌شود. توجه کنید که همهٔ !!{element}های
$\bigcup A$ پیش از مرحلهٔ $S$ در دسترس‌اند. اکنون بنا بر
\stagesacc{}، برای \emph{هر} مجموعه‌هایی که پیش از~$S$ تشکیل شده‌اند،
مجموعه‌ای تشکیل می‌شود که اعضایش دقیقاً همان‌ها هستند. به بیان دیگر،
هر \emph{گردایهٔ ممکن} از مجموعه‌های پیش‌تر در دسترس، در~$S$ وجود
خواهد داشت. اما مسلماً \emph{ممکن} است اشیایی را برگزینیم که بتوانند
مجموعه‌ای انتخاب برای~$A$ تشکیل دهند؛ این فقط زیرمجموعه‌ای کاملاً
مشخص از $\bigcup A$ است. پس چنان مجموعهٔ انتخابی، چنان‌که لازم بود،
وجود دارد.

خب، این کوششی \emph{بسیار} سریع برای توجیه انتخاب بر مبنایی درونی
بود. اما برای پیگیری بیشتر این اندیشه باید شرح زیبای پاتر
(\citeyear[\S14.8]{Potter2004}) را بخوانید.

\end{document}
